\documentclass[12pt,letterpaper]{article}

% =========================
% PAQUETES
% =========================
\usepackage[spanish, es-tabla]{babel}
\usepackage[utf8]{inputenc}
\usepackage[T1]{fontenc}
\usepackage{lmodern}
\usepackage{geometry}
\geometry{margin=2.5cm}
\usepackage{setspace}
\onehalfspacing
\usepackage{parskip}
\usepackage{graphicx}
\usepackage{array}
\usepackage{longtable}
\usepackage{booktabs}
\usepackage{tabularx}
\usepackage{multirow}
\usepackage{enumitem}
\usepackage{float}
\usepackage{titlesec}
\usepackage{xcolor}
\usepackage{hyperref}
\usepackage{fancyhdr}
\usepackage{lastpage}
\usepackage{amsmath}
\usepackage{pdflscape}
\usepackage{ragged2e}
\usepackage{caption}
\usepackage{colortbl}

% =========================
% CONFIGURACION
% =========================
\hypersetup{
    colorlinks=true,
    linkcolor=black,
    urlcolor=black,
    citecolor=black,
    pdftitle={Informe de Impacto Social, Territorial, Ambiental y Estratégico - Pista de Atletismo Rapa Nui},
    pdfauthor={Lucas Nervi},
    pdfsubject={Perfil de proyecto de inversión pública},
    pdfkeywords={Rapa Nui, pista de atletismo, inversión pública, impacto social, infraestructura deportiva}
}

\pagestyle{fancy}
\fancyhf{}
\fancyhead[L]{Informe de Impacto Social, Territorial, Ambiental y Estratégico}
\fancyhead[R]{Pista de Atletismo en Rapa Nui}
\fancyfoot[C]{\thepage\ de \pageref{LastPage}}

\titleformat{\section}{\large\bfseries}{\thesection.}{0.5em}{}
\titleformat{\subsection}{\normalsize\bfseries}{\thesubsection.}{0.5em}{}
\titleformat{\subsubsection}{\normalsize\itshape}{\thesubsubsection.}{0.5em}{}

\setlist[itemize]{leftmargin=1.3cm}
\setlist[enumerate]{leftmargin=1.3cm}

% =========================
% DATOS
% =========================
\title{\textbf{INFORME DE IMPACTO SOCIAL, TERRITORIAL, AMBIENTAL Y ESTRATÉGICO}\\
\large Proyecto: Construcción de Pista de Atletismo y Centro Atlético del Pacífico Sur en Rapa Nui}
\author{Documento base para presentación ante Municipalidad, Gobierno Regional, IND y organismos sectoriales}
\date{Marzo 2026}

\begin{document}

% =========================
% PORTADA
% =========================
\begin{titlepage}
    \centering
    \vspace*{2cm}
    {\Huge \textbf{INFORME DE IMPACTO SOCIAL, TERRITORIAL, AMBIENTAL Y ESTRATÉGICO}\par}
    \vspace{0.8cm}
    {\Large \textbf{Proyecto: Construcción de Pista de Atletismo y Centro Atlético del Pacífico Sur en Rapa Nui}\par}
    \vspace{1.5cm}
    {\large Documento de perfil técnico para formulación de iniciativa de inversión pública\par}
    \vspace{0.7cm}
    {\large Municipalidad de Rapa Nui / Gobierno Regional de Valparaíso / Instituto Nacional del Deporte\par}
    \vfill
    {\large Preparado como base de postulación y justificación institucional\par}
    \vspace{0.4cm}
    {\large Marzo 2026\par}
\end{titlepage}

\tableofcontents
\newpage

% =========================
% RESUMEN EJECUTIVO
% =========================
\section{Resumen Ejecutivo}

El presente informe desarrolla la justificación social, territorial, ambiental, económica y estratégica para la construcción de una pista de atletismo en la comuna de Isla de Pascua, Rapa Nui, concebida no sólo como una obra deportiva aislada, sino como el primer paso para la consolidación de un \textbf{Centro Atlético del Pacífico Sur}. La propuesta se formula bajo una lógica coherente con los criterios utilizados por el Sistema Nacional de Inversiones de Chile, particularmente en lo relativo a identificación del problema, población beneficiaria, análisis de alternativas, externalidades positivas, sostenibilidad y pertinencia territorial.

Rapa Nui constituye un territorio de excepcional singularidad geográfica, cultural y patrimonial. Su condición de aislamiento extremo incrementa los costos de provisión de infraestructura pública, pero al mismo tiempo vuelve más valiosa cada inversión que mejora calidad de vida, acceso a oportunidades y cohesión social. La isla presenta una necesidad estructural de fortalecer espacios que promuevan salud, desarrollo humano, formación de jóvenes y articulación comunitaria. En ese marco, el deporte es una herramienta de política pública especialmente eficaz, por su capacidad de incidir simultáneamente en salud, educación, convivencia, identidad y proyección territorial.

Actualmente, la comuna no dispone de infraestructura atlética especializada que permita entrenamiento sistemático, desarrollo escolar de pruebas atléticas, preparación de deportistas locales o realización de eventos deportivos de mediana escala. Esta carencia reduce las oportunidades de acceso efectivo a la práctica deportiva en condiciones adecuadas, obliga a adaptar espacios no diseñados para atletismo y limita el potencial de Rapa Nui para proyectarse como territorio de deporte, bienestar y turismo activo.

El proyecto propuesto considera la habilitación de un recinto con pista atlética de 400 metros, zonas de salto y lanzamiento, área de calentamiento, iluminación, servicios básicos, cierres, accesibilidad universal, bodegas y espacios de apoyo multipropósito. La infraestructura se concibe con una lógica de uso intensivo por parte de establecimientos educacionales, clubes, programas municipales, usuarios recreativos, deportistas de rendimiento y actividades comunitarias. En términos estratégicos, el proyecto puede escalarse gradualmente hacia un polo deportivo insular que conecte deporte, turismo responsable, identidad local y cooperación con otras islas del Pacífico.

La presente iniciativa presenta una condición especialmente favorable de madurez temprana, al contar ya con un \textbf{emplazamiento preliminar identificado en Fundo Vaitea} y con \textbf{referencias iniciales de costo asociadas a Resinsa} para orientar su dimensionamiento, lo que permite pasar desde una aspiración general a una propuesta concreta de infraestructura pública para Rapa Nui.

Desde la perspectiva social, la iniciativa generaría beneficios significativos en salud pública, prevención del sedentarismo, fortalecimiento de trayectorias juveniles, inclusión, equidad territorial y uso positivo del tiempo libre. Desde el punto de vista económico, contribuiría a diversificar la oferta de actividades de la isla, favorecer eventos y concentraciones deportivas, activar servicios asociados y reforzar la imagen de Rapa Nui como destino singular. En materia ambiental y patrimonial, el proyecto requiere una localización y diseño cuidadosos, compatibles con la sensibilidad arqueológica, paisajística y cultural del territorio, así como con los mecanismos de participación pertinentes.

Por todo lo anterior, la construcción de la pista de atletismo se presenta como una inversión pública de alto valor social, con impactos positivos multidimensionales, viabilidad política y gran potencia narrativa para ser presentada como una iniciativa emblemática de desarrollo para Rapa Nui.

% =========================
% MARCO DEL PROYECTO
% =========================
\section{Objeto del Informe y Alcance}

\subsection{Objeto del informe}

El presente documento tiene por finalidad servir como base técnico-argumental para la formulación y presentación de una iniciativa de inversión pública destinada a financiar el diseño, construcción y puesta en marcha de una pista de atletismo en Rapa Nui. Su propósito es entregar una estructura robusta de justificación para ser utilizada ante:

\begin{itemize}
    \item Municipalidad de Rapa Nui;
    \item Gobierno Regional de Valparaíso;
    \item Instituto Nacional del Deporte;
    \item Ministerio de Desarrollo Social y Familia;
    \item Ministerios sectoriales vinculados a infraestructura, patrimonio, deporte y desarrollo territorial;
    \item potenciales instituciones colaboradoras nacionales e internacionales.
\end{itemize}

\subsection{Alcance}

Este informe tiene carácter de \textbf{perfil estratégico avanzado}. No reemplaza los estudios de prefactibilidad, diseño de ingeniería, levantamiento topográfico, informe arqueológico, modelación financiera detallada ni tramitación ambiental o sectorial específica, pero sí entrega una estructura sólida para justificar la necesidad pública de la obra y orientar su formulación.

\subsection{Enfoque metodológico}

El análisis se estructura sobre cinco dimensiones complementarias:

\begin{enumerate}
    \item \textbf{Dimensión social:} beneficios sobre salud, desarrollo humano e integración.
    \item \textbf{Dimensión territorial:} reducción de brechas y equidad de acceso en territorio aislado.
    \item \textbf{Dimensión ambiental y patrimonial:} compatibilidad con el contexto ecológico y cultural.
    \item \textbf{Dimensión económica:} externalidades positivas, turismo deportivo y dinamización local.
    \item \textbf{Dimensión estratégica:} capacidad del proyecto de transformarse en infraestructura emblemática.
\end{enumerate}

% =========================
% CONTEXTO TERRITORIAL
% =========================
\section{Contexto Territorial de Rapa Nui}

\subsection{Condición geográfica y aislamiento}

Rapa Nui posee una condición geográfica única dentro del territorio chileno. Su localización en el Pacífico sur oriental determina un régimen de abastecimiento, movilidad y provisión de servicios altamente condicionado por la distancia con el continente. Esta singularidad impone costos logísticos altos, restricciones de escala y fuertes dependencias de planificación previa. Sin embargo, también entrega una identidad territorial extraordinariamente valiosa para iniciativas que vinculan patrimonio, bienestar y proyección internacional.

\subsection{Dimensión demográfica}

La comuna de Isla de Pascua registra una superficie aproximada de 163,6 km\textsuperscript{2}. En el Censo 2017 se registró una población de 7.750 habitantes, con concentración principal en Hanga Roa. En términos de infraestructura pública, esto configura una escala comunal donde una obra correctamente diseñada puede generar un impacto proporcionalmente muy alto sobre la vida cotidiana del territorio.

\subsection{Centralidad de Hanga Roa}

La concentración de población y servicios en Hanga Roa facilita pensar la pista de atletismo como un equipamiento de acceso transversal para estudiantes, clubes, usuarios recreativos, adultos mayores, programas de salud y actividades comunitarias. A diferencia de proyectos dispersos, una infraestructura de este tipo puede operar como nodo articulador, con una tasa de utilización relevante a lo largo de todo el año.

\subsection{Valor cultural y patrimonial}

Rapa Nui no es un territorio cualquiera. Se trata de un espacio de relevancia cultural global, con densidad patrimonial, paisaje arqueológico sensible e identidad indígena viva. Por ello, cualquier infraestructura pública debe ser concebida bajo criterios de respeto territorial, participación temprana, localización responsable y diseño compatible con la memoria, la visualidad y los usos del suelo de la isla.

% =========================
% PROBLEMA PUBLICO
% =========================
\section{Identificación del Problema Público}

\subsection{Problema central}

La comuna de Rapa Nui presenta un \textbf{déficit de infraestructura deportiva especializada para la práctica del atletismo}, lo que limita el acceso regular, seguro y de calidad a una disciplina base para el desarrollo físico, educativo y comunitario de la población.

\subsection{Manifestaciones del problema}

El problema descrito se expresa en diversas situaciones observables:

\begin{itemize}
    \item ausencia de una pista reglamentaria o semirreglamentaria para entrenamiento;
    \item uso de espacios no diseñados para pruebas atléticas;
    \item limitaciones para actividades escolares y formativas;
    \item menor proyección de talentos deportivos locales;
    \item baja capacidad para organizar competencias, encuentros o concentraciones;
    \item pérdida de oportunidades para vincular deporte con turismo, salud y comunidad.
\end{itemize}

\subsection{Causas estructurales}

Entre las principales causas se identifican las siguientes:

\begin{itemize}
    \item rezago histórico en infraestructura deportiva especializada;
    \item alta complejidad logística y costo de materiales en territorio insular;
    \item priorización de otras necesidades públicas urgentes;
    \item ausencia de un relato estratégico que conecte la pista con beneficios multisectoriales;
    \item falta de un proyecto técnicamente articulado y políticamente escalable.
\end{itemize}

\subsection{Efectos sociales del problema}

La carencia de infraestructura produce efectos de corto y largo plazo:

\begin{itemize}
    \item menor acceso a actividad física estructurada;
    \item menor adherencia deportiva de niños y jóvenes;
    \item pérdida de oportunidades de prevención en salud;
    \item debilitamiento de trayectorias de formación deportiva;
    \item subutilización del deporte como herramienta de integración.
\end{itemize}

% =========================
% JUSTIFICACION ESTRATEGICA
% =========================
\section{Justificación Estratégica del Proyecto}

\subsection{Más que una pista: una pieza de infraestructura estructurante}

La principal fortaleza de este proyecto es que no debe ser comunicado como una obra aislada, sino como una \textbf{infraestructura estructurante} para el desarrollo social, deportivo y territorial de Rapa Nui. Una pista de atletismo correctamente formulada no es sólo una superficie de entrenamiento: es una plataforma para salud pública, escuela, recreación, disciplina, identidad, actividades comunitarias y posicionamiento territorial.

\subsection{Equidad territorial}

El acceso a infraestructura pública de calidad no debe depender exclusivamente de la cercanía con grandes centros urbanos. En territorios insulares y alejados, cada obra pública tiene un efecto redistributivo más intenso, ya que corrige brechas históricas en acceso a equipamiento, servicios y oportunidades de desarrollo personal.

\subsection{Inversión con retornos múltiples}

El proyecto posee retornos de naturaleza múltiple:

\begin{itemize}
    \item \textbf{retorno social:} mejora de bienestar, salud y convivencia;
    \item \textbf{retorno educativo:} fortalecimiento de trayectorias escolares y deportivas;
    \item \textbf{retorno económico:} eventos, turismo y servicios asociados;
    \item \textbf{retorno simbólico:} señal potente de inversión pública en un territorio singular;
    \item \textbf{retorno político-institucional:} proyecto visible, defendible y alineado con equidad territorial.
\end{itemize}

\subsection{Narrativa de alto valor para postulación}

La narrativa recomendada para la postulación debe ser:

\begin{quote}
\textit{“La pista de atletismo de Rapa Nui no sólo responde a una necesidad deportiva. Se trata de una infraestructura pública habilitante para la salud, la formación de jóvenes, la cohesión comunitaria y la proyección del territorio como referente insular de deporte, bienestar y encuentro en el Pacífico Sur.”}
\end{quote}

\subsection{Ventaja habilitante del terreno disponible}

Un elemento especialmente favorable para la maduración de esta iniciativa es la existencia de un terreno propuesto en el sector de Fundo Vaitea, sobre terrenos donados para efectos del proyecto. En el contexto de la inversión pública, la disponibilidad preliminar de suelo representa una ventaja sustantiva, ya que reduce una de las principales barreras de entrada de proyectos de infraestructura: la incertidumbre sobre localización.

Contar con un emplazamiento posible desde etapas tempranas permite ordenar de mejor forma los estudios de factibilidad, anticipar compatibilidades territoriales, iniciar conversaciones institucionales con mayor claridad y proyectar el desarrollo del recinto bajo una lógica de crecimiento por etapas.

% =========================
% OBJETIVOS
% =========================
\section{Objetivos del Proyecto}

\subsection{Objetivo general}

Construir una infraestructura atlética pública en Rapa Nui que permita ampliar el acceso a la actividad física y al deporte, reducir brechas territoriales en equipamiento especializado y generar impactos positivos en salud, educación, cohesión social y desarrollo local.

\subsection{Objetivos específicos}

\begin{enumerate}
    \item Habilitar un recinto apto para entrenamiento atlético, uso escolar y actividades recreativas.
    \item Mejorar las condiciones de práctica deportiva para niños, jóvenes y adultos.
    \item Fortalecer la oferta pública de infraestructura deportiva de la comuna.
    \item Generar una base física para programas de desarrollo atlético y eventos deportivos.
    \item Posicionar a Rapa Nui como territorio con vocación para deporte, bienestar y encuentro insular.
\end{enumerate}

% =========================
% DESCRIPCION DEL PROYECTO
% =========================
\section{Descripción General del Proyecto}

\subsection{Componentes principales}

Se propone un complejo atlético de escala comunal, compuesto al menos por:

\begin{itemize}
    \item pista de atletismo de 400 metros;
    \item recta de velocidad y calles demarcadas;
    \item zonas de salto horizontal y vertical;
    \item zonas de lanzamiento con resguardos de seguridad;
    \item área de calentamiento y acondicionamiento físico;
    \item iluminación básica o intermedia según etapa;
    \item bodegas, servicios higiénicos y camarines;
    \item cierres perimetrales y control de acceso;
    \item accesibilidad universal;
    \item paisajismo y tratamiento de bordes.
\end{itemize}

\subsection{Escalabilidad por etapas}

Una ventaja central del proyecto es su posibilidad de implementación por fases:

\begin{enumerate}
    \item \textbf{Etapa 1:} habilitación base del terreno, pista y servicios mínimos;
    \item \textbf{Etapa 2:} áreas complementarias, iluminación y equipamiento ampliado;
    \item \textbf{Etapa 3:} consolidación como centro deportivo insular con programación anual.
\end{enumerate}

\subsection{Tipología de uso}

El recinto permitiría usos múltiples:

\begin{itemize}
    \item educación física;
    \item talleres municipales;
    \item deporte recreativo;
    \item entrenamiento competitivo;
    \item encuentros inter escolares;
    \item actividades de salud y envejecimiento activo;
    \item eventos comunitarios compatibles.
\end{itemize}

\subsection{Emplazamiento preliminar}

El emplazamiento preliminar considerado para el desarrollo del proyecto corresponde a un terreno ubicado en el sector de \textbf{Fundo Vaitea}, en Rapa Nui. De acuerdo con los antecedentes disponibles a la fecha, se trata de terrenos donados para efectos de materializar la iniciativa, condición que fortalece la viabilidad inicial del proyecto al disponer de una alternativa concreta de localización.

Sin perjuicio de lo anterior, la materialización definitiva del proyecto en dicho emplazamiento deberá quedar sujeta a los procesos de validación jurídica, factibilidad técnica, compatibilidad normativa, revisión patrimonial y evaluación territorial que correspondan. En este sentido, el terreno propuesto constituye una base favorable para la formulación de la iniciativa, pero no exime la necesidad de desarrollar los estudios previos requeridos para confirmar su idoneidad.

Desde una perspectiva estratégica, la localización en Fundo Vaitea presenta valor territorial por su asociación con una zona de gran escala y potencial de proyección, lo que permite pensar la obra no sólo como una pista aislada, sino como parte de una futura infraestructura deportiva de mayor alcance para la comuna.

% =========================
% POBLACION BENEFICIARIA
% =========================
\section{Población Objetivo y Beneficiarios}

\subsection{Beneficiarios directos}

Se considera como beneficiarios directos:

\begin{itemize}
    \item estudiantes de establecimientos educacionales de la comuna;
    \item usuarios de programas deportivos municipales;
    \item deportistas federados o en formación;
    \item clubes, agrupaciones y talleres de atletismo o acondicionamiento físico;
    \item usuarios recreativos regulares.
\end{itemize}

\subsection{Beneficiarios indirectos}

Como beneficiarios indirectos se identifican:

\begin{itemize}
    \item familias y redes de apoyo de niños y jóvenes participantes;
    \item servicios locales de educación, salud y turismo;
    \item comercio y emprendimientos asociados a eventos;
    \item instituciones públicas que requieren infraestructura multipropósito;
    \item visitantes y deportistas externos en programas de intercambio o concentración.
\end{itemize}

\subsection{Grupos prioritarios}

El proyecto es especialmente relevante para:

\begin{itemize}
    \item niñez y adolescencia;
    \item mujeres y jóvenes con acceso limitado a infraestructura especializada;
    \item personas mayores activas;
    \item usuarios con enfoque de inclusión;
    \item comunidad local interesada en hábitos de vida saludables.
\end{itemize}

% =========================
% IMPACTO SOCIAL
% =========================
\section{Impacto Social Esperado}

\subsection{Salud pública y prevención}

El impacto social más evidente de una pista de atletismo es su contribución a la actividad física regular. Esta dimensión tiene efectos sobre prevención del sedentarismo, salud cardiovascular, bienestar emocional, manejo del estrés y construcción de hábitos sostenibles. En territorios donde la escala poblacional es acotada, una sola infraestructura bien operada puede transformar significativamente los patrones de uso del tiempo libre.

\subsection{Desarrollo juvenil}

La infraestructura atlética ofrece una disciplina exigente, medible y formativa. El atletismo enseña constancia, autorregulación, tolerancia a la frustración, progresión por metas y valoración del esfuerzo. Para jóvenes de Rapa Nui, contar con un espacio propio de entrenamiento significa además reconocimiento, dignidad territorial y posibilidad de proyectar una identidad deportiva local.

\subsection{Cohesión comunitaria}

El proyecto fortalecería la vida comunitaria al ofrecer un espacio seguro, visible y legitimado de encuentro. La infraestructura deportiva, cuando es pública y programada, reduce fragmentación social, estimula participación y mejora la apropiación colectiva del espacio.

\subsection{Inclusión y acceso}

Una pista abierta a diversos programas puede acoger desde entrenamiento competitivo hasta caminatas guiadas, talleres escolares, pruebas comunales, actividades de salud y jornadas intergeneracionales. Esta diversidad de usos mejora la relación costo-impacto de la inversión y eleva su legitimidad social.

% =========================
% IMPACTO EDUCATIVO Y DEPORTIVO
% =========================
\section{Impacto Educativo y Deportivo}

\subsection{Fortalecimiento del sistema escolar}

Los establecimientos educacionales de la comuna podrían utilizar la pista para clases de educación física, pruebas atléticas, jornadas recreativas y torneos internos. Esto elevaría la calidad del trabajo pedagógico, permitiría mediciones más adecuadas y fortalecería la cultura deportiva desde edades tempranas.

\subsection{Detección y desarrollo de talentos}

El atletismo requiere continuidad, medición y espacio técnico. La existencia de una pista en la isla permitiría detectar y acompañar talentos deportivos sin que el primer paso formativo dependa necesariamente del traslado al continente.

\subsection{Base para otros deportes}

El atletismo es una disciplina base. Mejorar velocidad, resistencia, coordinación, fuerza y técnica de carrera beneficia transversalmente a múltiples deportes. Por ello, la pista tendría un efecto positivo sobre el ecosistema deportivo completo de Rapa Nui.

\subsection{Programación deportiva anual}

La inversión se vuelve más eficiente si desde su diseño se proyecta una malla mínima de uso:

\begin{itemize}
    \item bloques escolares en horario diurno;
    \item talleres formativos por la tarde;
    \item franjas de uso libre controlado;
    \item programas de salud y acondicionamiento para adultos;
    \item calendario anual de encuentros y actividades especiales.
\end{itemize}

% =========================
% IMPACTO ECONOMICO
% =========================
\section{Impacto Económico y Encadenamientos Locales}

\subsection{Efectos económicos directos e indirectos}

Aunque la rentabilidad social de este proyecto no debe limitarse a ingresos monetarios, sí existen externalidades económicas positivas relevantes:

\begin{itemize}
    \item gasto local asociado a construcción y mantención;
    \item consumo en comercio y servicios durante actividades deportivas;
    \item contratación de monitores, apoyo técnico y personal de operación;
    \item fortalecimiento de emprendimientos ligados a deporte, bienestar y turismo activo.
\end{itemize}

\subsection{Aumento de capacidad para eventos}

Una pista de atletismo no sólo sirve para uso cotidiano. También amplía la capacidad de la comuna para organizar actividades que atraigan visitantes, instituciones, delegaciones y programas públicos itinerantes. Cada actividad bien diseñada moviliza alojamiento, alimentación, transporte local y servicios complementarios.

\subsection{Diversificación de la economía territorial}

Rapa Nui depende fuertemente del turismo. Por ello, toda infraestructura que contribuya a diversificar la experiencia del territorio sin tensionar su patrimonio tiene valor estratégico. El deporte, especialmente en formatos compatibles con bienestar y vida activa, puede complementar la oferta tradicional y contribuir a desestacionalizar parcialmente algunos flujos.

% =========================
% TURISMO DEPORTIVO
% =========================
\section{Turismo Deportivo y Posicionamiento Territorial}

\subsection{Concepto de turismo deportivo compatible con Rapa Nui}

No se propone masificar el territorio mediante grandes competencias de alto impacto, sino desarrollar una línea de \textbf{turismo deportivo responsable, de escala compatible y alto valor simbólico}. Esto puede incluir:

\begin{itemize}
    \item campamentos de entrenamiento;
    \item encuentros escolares o universitarios;
    \item clínicas deportivas;
    \item festivales de deporte y bienestar;
    \item programas que combinen actividad física, cultura y territorio.
\end{itemize}

\subsection{Ventajas narrativas de Rapa Nui}

Pocas localizaciones en el mundo poseen el poder simbólico de Rapa Nui. Un recinto deportivo bien diseñado puede insertarse en una narrativa potente:

\begin{quote}
\textit{Deporte en uno de los territorios habitados más aislados del planeta, con identidad cultural única, paisaje extraordinario y vocación de encuentro del Pacífico.}
\end{quote}

\subsection{Oferta complementaria}

La pista permitiría articular programas que combinen:

\begin{itemize}
    \item entrenamiento;
    \item educación patrimonial;
    \item actividades al aire libre;
    \item intercambios deportivos;
    \item promoción de hábitos saludables.
\end{itemize}

% =========================
% BENCHMARK
% =========================
\section{Benchmark Estratégico: Infraestructura Atlética en Islas y Territorios Aislados}

\subsection{Justificación del benchmark}

El valor del benchmark no radica en copiar mecánicamente otros casos, sino en demostrar que los territorios insulares exitosos invierten en equipamientos multifuncionales que cumplen simultáneamente fines sociales, deportivos, educativos y turísticos.

\subsection{Criterios comparativos sugeridos}

Para construir un benchmark institucional convincente, se recomienda comparar Rapa Nui con territorios insulares o remotos bajo criterios como:

\begin{itemize}
    \item aislamiento geográfico;
    \item tamaño poblacional;
    \item centralidad del turismo;
    \item sensibilidad ambiental y patrimonial;
    \item necesidad de infraestructura multipropósito;
    \item capacidad de albergar eventos de pequeña y mediana escala.
\end{itemize}

\subsection{Aprendizajes transferibles}

Los principales aprendizajes de proyectos exitosos en territorios aislados suelen ser:

\begin{enumerate}
    \item diseñar obras multipropósito, no monofuncionales;
    \item asegurar programación anual desde el inicio;
    \item integrar identidad local al relato del proyecto;
    \item dimensionar mantención y logística desde fase temprana;
    \item vincular infraestructura con educación y comunidad, no sólo con competencia.
\end{enumerate}

\subsection{Aplicación al caso Rapa Nui}

El caso de Rapa Nui destaca por poseer una combinación especialmente atractiva de:

\begin{itemize}
    \item identidad cultural única;
    \item escala comunal donde la infraestructura tiene alta visibilidad;
    \item capacidad de construir un proyecto emblemático;
    \item posibilidad de liderar una conversación sobre deporte en territorios aislados.
\end{itemize}

% =========================
% CENTRO PACIFICO SUR
% =========================
\section{Proyección Estratégica: Centro Atlético del Pacífico Sur}

\subsection{Cambio de escala conceptual}

La formulación recomendada no es \textit{``solo una pista''}. El enfoque estratégico consiste en presentar la obra como la primera infraestructura habilitante de un \textbf{Centro Atlético del Pacífico Sur}, concebido como polo de:

\begin{itemize}
    \item formación atlética;
    \item intercambio deportivo;
    \item deporte escolar insular;
    \item bienestar y vida activa;
    \item investigación aplicada a deporte en contextos remotos.
\end{itemize}

\subsection{Ventajas del concepto}

Este cambio de escala mejora significativamente el peso político del proyecto, porque:

\begin{itemize}
    \item conecta la obra con una visión de futuro;
    \item aumenta la capacidad de atraer apoyos interinstitucionales;
    \item permite construir un relato diferenciador frente a otras inversiones;
    \item evita que el proyecto sea percibido como una infraestructura secundaria.
\end{itemize}

\subsection{Líneas futuras posibles}

A mediano plazo, el Centro Atlético del Pacífico Sur podría incorporar:

\begin{itemize}
    \item sala multipropósito para preparación física;
    \item unidad básica de evaluación deportiva;
    \item espacios de educación y clínica deportiva;
    \item residencias cortas para concentraciones o pasantías;
    \item colaboración con universidades y programas de salud.
\end{itemize}

% =========================
% IMPACTO AMBIENTAL
% =========================
\section{Impacto Ambiental}

\subsection{Enfoque general}

La iniciativa presenta impactos ambientales previsibles y gestionables, concentrados principalmente en la fase de construcción. No obstante, dado el contexto territorial de Rapa Nui, la evaluación ambiental no debe limitarse a una lógica de cumplimiento mínimo, sino incorporar desde el diseño una cultura de bajo impacto, compatibilidad paisajística y manejo responsable del suelo.

\subsection{Impactos potenciales en fase de construcción}

\begin{itemize}
    \item movimiento de tierras;
    \item generación de material particulado;
    \item ruido temporal por faenas;
    \item generación de residuos de construcción;
    \item tránsito de maquinaria y acopio de materiales.
\end{itemize}

\subsection{Impactos potenciales en fase de operación}

\begin{itemize}
    \item consumo de agua para mantención de áreas complementarias;
    \item generación de residuos por uso del recinto;
    \item demanda eléctrica por iluminación;
    \item presión de uso sobre accesos y circulación.
\end{itemize}

\subsection{Lineamientos de mitigación}

Se recomienda incorporar desde la etapa de diseño:

\begin{itemize}
    \item plan de manejo de residuos;
    \item control de polvo y material particulado;
    \item medidas de eficiencia hídrica;
    \item soluciones de iluminación eficiente;
    \item tratamiento paisajístico con especies compatibles;
    \item diseño de drenaje adecuado;
    \item manual de operación sustentable del recinto.
\end{itemize}

% =========================
% PATRIMONIO Y CULTURA
% =========================
\section{Resguardo Patrimonial, Cultural y Paisajístico}

\subsection{Principio rector}

En Rapa Nui, el patrimonio no es un componente secundario del territorio: es un eje estructural de cualquier intervención pública. Por ello, la viabilidad del proyecto dependerá en gran medida de su correcta localización, del tratamiento de bordes, del lenguaje arquitectónico y de la existencia de una narrativa de respeto y beneficio comunitario.

\subsection{Criterios esenciales}

El proyecto deberá incorporar, como mínimo, los siguientes criterios:

\begin{enumerate}
    \item localización fuera de áreas sensibles o con interferencia arqueológica;
    \item evaluación temprana de riesgos patrimoniales;
    \item diseño de baja agresividad visual;
    \item integración con topografía y materialidad del entorno;
    \item coherencia con la identidad cultural local.
\end{enumerate}

\subsection{Relación con el paisaje cultural}

Una infraestructura deportiva en Rapa Nui debe evitar competir visual o simbólicamente con el patrimonio. Su virtud debe ser la \textbf{discreción bien diseñada}: una obra útil, contemporánea y sobria, al servicio de la comunidad.

% =========================
% PARTICIPACION INDIGENA
% =========================
\section{Participación Comunitaria y Pertinencia Indígena}

\subsection{Participación como condición de legitimidad}

La participación comunitaria no debe ser entendida únicamente como un requisito formal, sino como una condición esencial de legitimidad política y social. Mientras más temprana y real sea la conversación con la comunidad, mayor será la solidez del proyecto en sus etapas posteriores.

\subsection{Actores que debiesen ser considerados}

\begin{itemize}
    \item Municipalidad de Rapa Nui;
    \item organizaciones deportivas locales;
    \item establecimientos educacionales;
    \item representantes comunitarios;
    \item actores vinculados a patrimonio y cultura;
    \item servicios públicos competentes;
    \item eventuales usuarios prioritarios del recinto.
\end{itemize}

\subsection{Temas clave a dialogar}

\begin{itemize}
    \item ubicación;
    \item usos prioritarios;
    \item horarios y sistema de administración;
    \item resguardo cultural y paisajístico;
    \item beneficios directos para la población local;
    \item programación anual;
    \item escalamiento futuro del proyecto.
\end{itemize}

\subsection{Resultado esperado}

La participación bien conducida permite transformar el proyecto desde una infraestructura “para” la comunidad a una infraestructura “con” la comunidad.

% =========================
% ALTERNATIVAS
% =========================
\section{Análisis de Alternativas}

\subsection{Alternativa 1: Mantener situación actual}

Consiste en continuar utilizando infraestructura no especializada. Esta alternativa presenta el menor costo inmediato, pero perpetúa la brecha, limita el desarrollo deportivo y mantiene el déficit de calidad en la oferta pública.

\subsection{Alternativa 2: Mejoramiento parcial de espacios existentes}

Implica adaptar o mejorar espacios ya disponibles. Puede tener menor costo inicial, pero usualmente produce rendimientos limitados, usos incompatibles y una vida útil funcional menor en términos de atletismo.

\subsection{Alternativa 3: Construcción de pista de atletismo con enfoque multipropósito}

Es la alternativa recomendada. Permite resolver de manera estructural el problema, aumentar la capacidad comunal, generar mayor impacto social y construir una obra escalable.

\subsection{Síntesis comparativa}

\begin{table}[H]
\centering
\caption{Comparación cualitativa de alternativas}
\begin{tabularx}{\textwidth}{>{\RaggedRight\arraybackslash}p{3.5cm}XXX}
\toprule
\textbf{Criterio} & \textbf{Situación actual} & \textbf{Mejoramiento parcial} & \textbf{Pista multipropósito} \\
\midrule
Resolución del problema & Baja & Media & Alta \\
Impacto social & Bajo & Medio & Alto \\
Escalabilidad & Nula & Baja & Alta \\
Visibilidad política & Baja & Media & Alta \\
Potencial deportivo & Bajo & Medio & Alto \\
Aporte territorial & Bajo & Medio & Alto \\
\bottomrule
\end{tabularx}
\end{table}

% =========================
% MODELO DE GESTION
% =========================
\section{Modelo de Gestión y Gobernanza}

\subsection{Principio de operación}

Una infraestructura de este tipo requiere un modelo de gestión claro para evitar subutilización. Se recomienda una gobernanza municipal con articulación intersectorial y calendario anual de uso.

\subsection{Esquema sugerido}

\begin{itemize}
    \item \textbf{Administrador principal:} Municipalidad de Rapa Nui.
    \item \textbf{Aliados programáticos:} establecimientos educacionales, IND, organizaciones deportivas, salud municipal.
    \item \textbf{Consejo consultivo local:} instancia liviana para seguimiento, priorización y mejora.
\end{itemize}

\subsection{Principios de uso}

\begin{enumerate}
    \item prioridad para población local;
    \item equilibrio entre uso libre y uso programado;
    \item resguardo del recinto y sostenibilidad operativa;
    \item enfoque inclusivo;
    \item programación visible y transparente.
\end{enumerate}

% =========================
% SOSTENIBILIDAD OPERATIVA
% =========================
\section{Sostenibilidad Operativa y Mantención}

\subsection{Condición crítica de éxito}

La principal amenaza para este tipo de proyectos no suele ser su construcción, sino su operación. Por ello, la sostenibilidad debe abordarse desde la fase de formulación.

\subsection{Factores críticos}

\begin{itemize}
    \item presupuesto anual de mantención;
    \item provisión de insumos y reposiciones;
    \item personal mínimo de operación;
    \item reglamento de uso;
    \item vigilancia y control de accesos;
    \item programación continua que legitime la inversión.
\end{itemize}

\subsection{Recomendaciones}

\begin{itemize}
    \item dimensionar costos de ciclo de vida;
    \item prever mantención preventiva;
    \item establecer alianzas para uso intensivo;
    \item asignar responsable operativo desde antes de inaugurar.
\end{itemize}

% =========================
% RIESGOS
% =========================
\section{Riesgos del Proyecto y Estrategias de Mitigación}

\begin{table}[H]
\centering
\caption{Matriz preliminar de riesgos}
\begin{tabularx}{\textwidth}{>{\RaggedRight\arraybackslash}p{4cm}>{\RaggedRight\arraybackslash}p{3cm}X}
\toprule
\textbf{Riesgo} & \textbf{Nivel} & \textbf{Medida de mitigación} \\
\midrule
Sobrecosto logístico de materiales & Alto & Planificación temprana de abastecimiento, compras consolidadas y diseño constructivo adaptado al territorio. \\
Interferencia patrimonial o arqueológica & Alto & Estudios previos, selección cuidadosa del terreno y coordinación temprana con organismos competentes. \\
Subutilización del recinto & Medio & Modelo de gestión con programación anual obligatoria y alianzas institucionales. \\
Deterioro prematuro por mantención insuficiente & Medio & Presupuesto de operación, mantención preventiva y responsable dedicado. \\
Conflictos por priorización de uso & Medio & Reglamento público de uso y gobernanza clara. \\
Resistencia comunitaria por falta de participación & Alto & Participación temprana, comunicación transparente y pertinencia cultural. \\
\bottomrule
\end{tabularx}
\end{table}

\subsection{Riesgo jurídico y de regularización del terreno}

Si bien la disponibilidad preliminar de terrenos donados en Fundo Vaitea constituye una fortaleza importante del proyecto, será indispensable verificar y formalizar su situación jurídica, condiciones de uso, eventuales restricciones normativas y compatibilidad con el destino propuesto. La mitigación de este riesgo requiere revisión temprana de títulos, actos de donación, condiciones de administración y factibilidad institucional para su incorporación efectiva a una iniciativa pública de inversión.

% =========================
% INDICADORES
% =========================
\section{Indicadores de Resultado e Impacto}

\subsection{Indicadores de corto plazo}

\begin{itemize}
    \item número de usuarios semanales;
    \item número de establecimientos que utilizan el recinto;
    \item horas mensuales de programación formal;
    \item número de talleres o escuelas deportivas activas.
\end{itemize}

\subsection{Indicadores de mediano plazo}

\begin{itemize}
    \item aumento de participación deportiva comunal;
    \item número de eventos realizados por año;
    \item tasa de uso por grupos prioritarios;
    \item continuidad de jóvenes en programas deportivos.
\end{itemize}

\subsection{Indicadores estratégicos}

\begin{itemize}
    \item posicionamiento de Rapa Nui como sede de actividades deportivas de escala compatible;
    \item integración del recinto en planes comunales de salud, educación y deporte;
    \item grado de consolidación del Centro Atlético del Pacífico Sur.
\end{itemize}

% =========================
% PRESUPUESTO REFERENCIAL
% =========================
\section{Estructura Presupuestaria Referencial}

\subsection{Criterio general}

A la fecha de elaboración del presente informe, se cuenta con una \textbf{referencia preliminar de costos} asociada a la empresa \textbf{Resinsa}, vinculada a Werner Schmidt Tuki, antecedente que constituye una base inicial útil para la estimación temprana del proyecto en etapa de perfil.

Este antecedente debe entenderse como una \textbf{referencia presupuestaria preliminar no definitiva}, cuya utilidad principal es orientar la magnitud esperada de inversión y fortalecer la discusión institucional sobre viabilidad, escalabilidad y fuentes de financiamiento. El costo definitivo del proyecto deberá establecerse en etapas posteriores mediante desarrollo de diseño, especificaciones técnicas, cubicaciones, presupuesto detallado y eventuales procesos de cotización o licitación que correspondan conforme a la normativa aplicable.

\subsection{Uso del costo referencial en etapa de perfil}

En la fase actual, la referencia de costo entregada por Resinsa permite:

\begin{itemize}
    \item dimensionar preliminarmente la inversión requerida;
    \item sustentar la conversación con autoridades municipales, regionales y sectoriales;
    \item comparar escenarios de alcance base y alcance ampliado;
    \item justificar la necesidad de avanzar hacia estudios de prefactibilidad y diseño.
\end{itemize}

\subsection{Alcance sugerido de la estimación referencial}

Se recomienda que la referencia preliminar de costo sea ordenada, para efectos de presentación institucional, en los siguientes componentes:

\begin{itemize}
    \item habilitación y preparación del terreno;
    \item movimiento de tierra y nivelación;
    \item construcción de pista y base estructural;
    \item obras complementarias para saltos y lanzamientos;
    \item cierres perimetrales y accesos;
    \item iluminación y urbanización básica;
    \item servicios de apoyo, bodegas y camarines;
    \item equipamiento deportivo;
    \item paisajismo e integración al entorno;
    \item gastos generales, transporte y contingencias logísticas asociadas a la condición insular.
\end{itemize}

\subsection{Recomendación metodológica}

Para efectos de postulación a fondos públicos, se recomienda trabajar con dos escenarios presupuestarios:

\begin{enumerate}
    \item \textbf{Escenario base:} construcción de pista y operación esencial del recinto.
    \item \textbf{Escenario ampliado:} construcción de pista con componentes complementarios que fortalezcan la visión de Centro Atlético del Pacífico Sur.
\end{enumerate}

Esta distinción permite una mejor conversación con los organismos financiadores, facilitando la priorización de una primera etapa financiable sin renunciar a una visión estratégica de largo plazo.

% =========================
% HOJA DE RUTA
% =========================
\section{Hoja de Ruta Recomendada para Postulación}

\subsection{Fase 1: Validación política e institucional}

\begin{enumerate}
    \item Aprobación conceptual municipal.
    \item Definición del relato estratégico oficial.
    \item Priorización del proyecto en cartera comunal/regional.
\end{enumerate}

\subsection{Fase 2: Validación territorial}

\begin{enumerate}
    \item identificación preliminar de terreno;
    \item conversación temprana con actores locales;
    \item evaluación inicial de compatibilidades territoriales.
\end{enumerate}

\subsection{Fase 3: Estudios habilitantes}

\begin{enumerate}
    \item topografía;
    \item mecánica de suelos;
    \item revisión patrimonial/arquelógica;
    \item perfil técnico-económico;
    \item estimación de costos de ciclo de vida.
\end{enumerate}

\subsection{Fase 4: Formulación y postulación}

\begin{enumerate}
    \item ingreso de iniciativa;
    \item formulación bajo metodología sectorial correspondiente;
    \item validación presupuestaria;
    \item búsqueda de financiamiento.
\end{enumerate}

\subsection{Fase 5: Diseño, ejecución y activación}

\begin{enumerate}
    \item licitación de diseño y obras;
    \item construcción;
    \item reglamento y plan de operación;
    \item calendario de uso desde la inauguración.
\end{enumerate}

% =========================
% ARGUMENTARIO INSTITUCIONAL
% =========================
\section{Argumentario Institucional para Autoridades}

\subsection{Mensaje para municipalidad}

\textit{Este proyecto entrega a la comuna una infraestructura visible, útil y transversal, con capacidad de impacto real sobre niñez, juventud, salud y vida comunitaria.}

\subsection{Mensaje para Gobierno Regional}

\textit{La iniciativa reduce brechas territoriales en un territorio extremo, combina equidad con alto valor simbólico y puede transformarse en una obra emblemática de descentralización efectiva.}

\subsection{Mensaje para sector deporte}

\textit{La pista crea base física para programas formativos, deporte escolar, desarrollo atlético y articulación con otras disciplinas.}

\subsection{Mensaje para patrimonio y comunidad}

\textit{El proyecto puede ser una obra pública contemporánea, sobria y respetuosa, diseñada con participación local y compatible con la identidad territorial de Rapa Nui.}

% =========================
% CONCLUSION
% =========================
\section{Conclusiones}

La construcción de una pista de atletismo en Rapa Nui constituye una iniciativa pública de alta pertinencia y gran potencial transformador. Su valor no reside únicamente en resolver una carencia de infraestructura deportiva, sino en habilitar una plataforma concreta para salud, educación, cohesión comunitaria, equidad territorial y proyección estratégica del territorio.

El proyecto cuenta con una lógica especialmente potente para efectos de postulación: responde a una necesidad real, beneficia de forma transversal a la comunidad, tiene externalidades positivas múltiples, puede ejecutarse por etapas y posee un relato institucional altamente defendible. Adicionalmente, su escalamiento conceptual como \textbf{Centro Atlético del Pacífico Sur} permite dotarlo de una visión de largo plazo capaz de movilizar apoyos políticos y técnicos de mayor envergadura.

Adicionalmente, la existencia de un emplazamiento preliminar en Fundo Vaitea y de una referencia inicial de costos asociada a Resinsa mejora de manera significativa la madurez temprana de la iniciativa, al permitir avanzar desde una idea general hacia una propuesta territorialmente identificable y financieramente dimensionable.

La recomendación central de este informe es avanzar en su formulación oficial con un enfoque doble: por una parte, como infraestructura deportiva indispensable para cerrar una brecha concreta; por otra, como proyecto emblemático de inversión pública en un territorio singular de Chile. En esa combinación se encuentra su principal fuerza.

% =========================
% ANEXOS
% =========================
\newpage
\section*{Anexo A: Información mínima que conviene agregar antes de presentar}
\addcontentsline{toc}{section}{Anexo A: Información mínima que conviene agregar antes de presentar}

Antes de entregar este documento como versión final, se recomienda completar los siguientes elementos con información local validada:

\begin{itemize}
    \item terreno propuesto y justificación de localización;
    \item croquis o plano de emplazamiento;
    \item fotografías del área actual;
    \item carta de apoyo municipal o institucional;
    \item estimación preliminar de costos;
    \item número de establecimientos educacionales beneficiarios;
    \item programas deportivos actuales de la comuna;
    \item diagnóstico breve de uso esperado por segmento;
    \item ruta de permisos y compatibilidades patrimoniales.
\end{itemize}

\section*{Anexo B: Propuesta de cartera inicial de usos del recinto}
\addcontentsline{toc}{section}{Anexo B: Propuesta de cartera inicial de usos del recinto}

\begin{longtable}{p{5cm}p{9cm}}
\toprule
\textbf{Línea programática} & \textbf{Uso propuesto} \\
\midrule
Deporte escolar & Pruebas atléticas, talleres y encuentros inter escolares. \\
Formación deportiva & Escuela municipal de atletismo por categorías. \\
Salud comunitaria & Caminatas guiadas, acondicionamiento físico, programas para adultos. \\
Juventud & Talleres vespertinos de velocidad, resistencia y fuerza básica. \\
Comunidad & Corridas familiares, jornadas de actividad física, festivales deportivos. \\
Turismo deportivo & Clínicas, campamentos y programas de intercambio de escala compatible. \\
\bottomrule
\end{longtable}

\section*{Anexo C: Texto breve de presentación política}
\addcontentsline{toc}{section}{Anexo C: Texto breve de presentación política}

\textbf{Versión sugerida:}

La construcción de una pista de atletismo en Rapa Nui es una obra pública con impacto real y transversal. Fortalece salud, formación de jóvenes, vida comunitaria y equidad territorial en uno de los territorios más aislados de Chile. Además, proyecta a la isla como un referente insular de deporte y bienestar en el Pacífico Sur. No se trata sólo de una infraestructura deportiva: se trata de una inversión estratégica para la calidad de vida y el desarrollo futuro de la comunidad.

\section*{Anexo D: Antecedentes preliminares de formulación}
\addcontentsline{toc}{section}{Anexo D: Antecedentes preliminares de formulación}

Para efectos de la presente versión del informe, se consideran como antecedentes preliminares de formulación los siguientes elementos de trabajo:

\begin{itemize}
    \item identificación de un emplazamiento preliminar en el sector Fundo Vaitea;
    \item existencia de terrenos donados para fines asociados al proyecto;
    \item referencia preliminar de costos asociada a la empresa Resinsa, vinculada a Werner Schmidt Tuki, utilizada únicamente como base inicial de dimensionamiento de inversión.
\end{itemize}

Estos antecedentes deberán ser formalizados, validados y complementados en las etapas posteriores de prefactibilidad, diseño y formulación oficial de la iniciativa.

% =========================
% BIBLIOGRAFIA
% =========================
\newpage
\begin{thebibliography}{99}

\bibitem{SNIdeportes}
Sistema Nacional de Inversiones. \textit{Metodología Sector Deportes}. Gobierno de Chile. Disponible en:
\url{https://sni.gob.cl/storage/docs/Metodologi%CC%81a-Deportes-2016.pdf}

\bibitem{SEAguias}
Servicio de Evaluación Ambiental. \textit{Guías para la evaluación de impacto ambiental}. Disponible en:
\url{https://www.sea.gob.cl/guias-para-la-evaluacion-de-impacto-ambiental}

\bibitem{SEAgruposhumanos}
Servicio de Evaluación Ambiental. \textit{Guía para la predicción y evaluación de impactos sobre los sistemas de vida y costumbres de grupos humanos en el SEIA}. Disponible en:
\url{https://www.sea.gob.cl/documentacion/guias-y-criterios/guia-para-la-prediccion-y-evaluacion-de-impactos-sobre-los-sistemas}

\bibitem{ConsultaIndigena}
Ministerio del Medio Ambiente. \textit{Consulta a los Pueblos Indígenas}. Disponible en:
\url{https://consultaindigena.mma.gob.cl/}

\bibitem{UNESCO}
UNESCO World Heritage Centre. \textit{Rapa Nui National Park}. Disponible en:
\url{https://whc.unesco.org/en/list/715/}

\bibitem{CONAF}
Corporación Nacional Forestal. \textit{Parque Nacional Rapa Nui}. Disponible en:
\url{https://www.conaf.cl/parque_nacionales/parque-nacional-rapa-nui/}

\bibitem{BCN2024}
Biblioteca del Congreso Nacional de Chile. \textit{Reporte Comunal 2024: Isla de Pascua}. Disponible en:
\url{https://www.bcn.cl/siit/reportescomunales/comunas_v.html?anno=2024&idcom=5201}

\bibitem{BCN2020}
Biblioteca del Congreso Nacional de Chile. \textit{Reporte Comunal: Isla de Pascua, población Censo 2017 y proyección}. Disponible en:
\url{https://www.bcn.cl/siit/reportescomunales/comunas_v.html?anno=2020&idcom=5201}

\bibitem{PoliticaDeporte}
Instituto Nacional del Deporte. \textit{Política Nacional de Actividad Física y Deporte 2016-2025}. Disponible en:
\url{https://sigi.ind.cl/secciones/142}

\end{thebibliography}

\end{document}